\section{Results}
\label{sec:results}

% ~1400 words. ORDER MATTERS: architecture first, because everything
% downstream is conditional on which quantiser this is.
%
% Split the section explicitly into CONFIRMATORY and EXPLORATORY. See
% PREREGISTRATION.md. A reader must be able to tell which is which without
% reconstructing the project's history.

\subsection{Architecture identification}
\label{sec:arch}
% CONFIRMATORY. V0.4 was scored against five pre-specified outcomes
% {rounding, truncation, RPDF, TPDF, neither}.
%
% gyro16 == gyro20 >> 4 exactly. \vzeroFourTotal{} of \vzeroFourTotal{}
% discriminating samples, three ODR spanning 40x, zero ambiguous.
% This is bit-exact and needs no statistics at all -- which is why rule R6's
% likelihood machinery is not invoked (see PREREGISTRATION.md, deviation 4).
\todo{Write.}

\subsection{Added power against dither ratio}
% \nOdrPoints{} axis-measurements, two specimens, rho over \rhoDecades{}
% decades. Residual RMS \odrResidRMS{} against exact theory.
\todo{Write. \cref{fig:etarho}.}

\subsection{Added power against phase: the causal manipulation}
\label{sec:phasesweep}
% THE result. One specimen, one ODR, one configuration; the only thing that
% changes between records is a number written to a trim register.
% \nPhasePoints{} measurements, eta from \etaRangeLo{} to \etaRangeHi{},
% residual RMS \phaseResidRMS{} with ZERO free parameters.
%
% This is what answers Objection #14 part 2 (Objections v2.1 Z.2): phi is
% MANIPULATED, not merely observed.
\todo{Write.}

\begin{figure}[t]\centering
  \includegraphics[width=\linewidth]{fig7_phase_sweep.png}
  \caption{\todo{Controlled phase sweep. Curve is exact theory; nothing
  fitted.}}
  \label{fig:phasesweep}
\end{figure}

\subsection{An independent channel}
% Code occupancy, max discrepancy \histMaxDisc{}, nothing fitted. A
% DISTRIBUTIONAL statistic sharing no algebra with eta.
\todo{Write. Figure to supplementary: fig15\_code\_histograms.png}

\subsection{The Allan domain}
% The absence of a -1 slope, with both reference lines anchored ON the data
% so the only thing being compared is the slope.
\todo{Write.}

\begin{figure}[t]\centering
  \includegraphics[width=\linewidth]{fig10_allan_family.png}
  \caption{\todo{The quantisation term is on the $-1/2$ family. No
  IEEE-952 $Q$ term appears.}}
  \label{fig:allan}
\end{figure}
